\section{Posterior replicas and exact equal-label measures}
\label{sec:replicas}

The information estimate for a block has two parts. First, a general
identity converts information in one finite message into the loss of
dependence among several signals. Second, exact Gaussian observations give
an explicit law for signals sampled from the same posterior. This section
proves both statements. The subsequent sections estimate how much of that
dependence remains visible after another projection.

All probability spaces used below are standard Borel. Conditional laws mean
measurable regular conditional probabilities, and conditional relative
entropies include integration over the conditioning variable. We write
$\KL(P\Vert Q)$ for relative entropy, with value $+\infty$ when $P$ is
not absolutely continuous with respect to $Q$. Logarithms and entropies are
in nats.

\subsection{Information in a message and loss of correlation}

For a joint law $Q$ of $(S_1,\ldots,S_t)$ with common marginal $p$, its
\emph{total correlation} is
\[
 K=\KL(Q\Vert p^{\otimes t}).
\]
This quantity was introduced by Watanabe~\cite{Watanabe1960}; see
Austin~\cite[Sections 3.3 and 6]{Austin2020} for its measure-theoretic
form and information identities. Let $B$ be independent of the tuple and
let $U_i=\phi_B(S_i)$, where $(B,s)\mapsto\phi_B(s)$ is measurable.
Writing $p_B=(\phi_B)_\#p$, define the correlation visible after projection by
\[
 K'=\E_B\KL\bigl(Q_{U_1,\ldots,U_t\mid B}\Vert p_B^{\otimes t}\bigr).
\]

\begin{lemma}[Replica information inequality]\label{lem:replica-information}
Let $t\ge1$ and suppose $K<\infty$. Let $W$ be finite-valued and suppose
that every pair $(S_i,W)$ has the same law as $(S_1,W)$. In the conclusion
write $S=S_1$. If $B$ is independent of $(S_1,\ldots,S_t,W)$, then
\begin{equation}\label{eq:replica-information}
 t I(S;W\mid B,\phi_B(S))\le H(W)+K-K',\qquad 0\le K'\le K.
\end{equation}
The same statement holds in a conditional experiment when all of these
hypotheses hold there, including the matching pair laws and the independence
of $B$. Equivalently, the left side may use a separate representative of
the common signal--message pair law if it is coupled independently with $B$.
\end{lemma}

\begin{proof}
Let $K_W$ be the average total correlation conditional on $W$, computed
relative to the product of the conditional one-coordinate marginals.
Decomposing the relative entropy of $(S_1,\ldots,S_t,W)$ from
$p^{\otimes t}\otimes Q_W$ in the two orders gives
\begin{equation}\label{eq:replica-chain-before}
 tI(S;W)=I((S_i)_{i=1}^t;W)+K-K_W.
\end{equation}
The relative entropy being decomposed is at most $K+H(W)$, so every
nonnegative chain-rule term in this display is finite. Define $K'_W$ in
the same way after projection, conditionally on $(B,W)$. The corresponding
identity after projection is
\begin{equation}\label{eq:replica-chain-after}
 tI(\phi_B(S);W\mid B)
 =I((U_i)_{i=1}^t;W\mid B)+K'-K'_W.
\end{equation}
Coordinatewise data processing gives $K'\le K$ and $K'_W\le K_W$.
Independence of $B$ and the fact that $\phi_B(S)$ is determined by $(B,S)$
also give
\[
 I(S;W\mid B,\phi_B(S))=I(S;W)-I(\phi_B(S);W\mid B).
\]
Subtracting \eqref{eq:replica-chain-after} from
\eqref{eq:replica-chain-before} therefore gives the exact identity
\[
 tI(S;W\mid B,\phi_B(S))
 =I((S_i)_{i=1}^t;W\mid B,(U_i)_{i=1}^t)
   +(K-K')-(K_W-K'_W).
\]
The conditional tuple-information is at most $H(W)$, and
$K_W-K'_W\ge0$. This proves the inequality.
The remaining assertion is nonnegativity of relative entropy.
\end{proof}

The coupling used in this paper has precisely the pair-law property in the
lemma. Suppose that data $(A,Y)$ are generated from $S$, and that $W$ is
drawn from a kernel of the data, conditionally independently of $S$.
Given the same $(A,Y)$, draw $S_1,\ldots,S_t$ independently from the
posterior law of $S$, independently also of $W$. Conditioning on the data
shows that $(S_i,W)$ has the original signal--message law for every $i$.
We call these signals \emph{posterior replicas}. Equal one-signal marginals
alone would not be sufficient for Lemma~\ref{lem:replica-information}.
The invariance of the joint data--signal law under posterior resampling
is the Nishimori identity; see Lelarge and
Miolane~\cite[Proposition 16]{LelargeMiolane2019} for a general formulation.
When an independent $B$ is added to both experiments, equality of the pair
laws gives equality of the corresponding $(S_i,W,B)$ laws as well, so the
lemma's left side is the information in the original signal--message pair.

We shall use the following entropy test. If $Q\ll R$ are probabilities and
$F\ge0$ is measurable with $\E_R F\le1$, then
\begin{equation}\label{eq:entropy-test}
 \KL(Q\Vert R)\ge\E_Q\log F
\end{equation}
whenever the right side is defined. For a bounded positive test this follows
by tilting $R$ by $F/\E_R F$ and using nonnegativity of relative entropy;
bounded truncations give the stated form. This is the entropy variational
principle in the form used here; see
Dupuis and Mao~\cite[Equations (1.1)--(1.2)]{DupuisMao2022}.
In particular, if the relative entropy is finite, applying the bounded form
to truncations of $(1+F)/2$ and then monotone convergence yields
\begin{equation}\label{eq:entropy-positive-log}
 \E_Q\log(1+F)\le\KL(Q\Vert R)+\log2.
\end{equation}
This supplies positive-log integrability for the likelihood tests below.

\subsection{Projection fibers and inverse affine volumes}

Let $\sigma_d$ be uniform probability on $S^{d-1}$ and let $\gamma_{k,d}$
be the law of a $k\times d$ matrix with independent standard Gaussian
entries. For a subspace $F\subset\R^d$, let $\gamma_{k,d}^{F}$ denote the
law of a matrix whose rows are independent standard Gaussians in $F^\perp$.
For a tuple $\mathbf u=(u_1,\ldots,u_t)$ with $m=t-1$, set
\begin{equation}\label{eq:equal-label-geometry}
 D_{\mathbf u}=[u_2-u_1,\ldots,u_t-u_1],\qquad
 F_{\mathbf u}=\col(D_{\mathbf u}),\qquad
 J(\mathbf u)=\det(D_{\mathbf u}^{\mathsf T}D_{\mathbf u})^{1/2}.
\end{equation}
At full column rank $J$ is the product of the successive distances to
predecessor affine spans. For two distinct points, the difference
$v=u_2-u_1$ has $Av\sim N(0,\norm v^2I_k)$, whose density at zero is
$(2\pi)^{-k/2}\norm v^{-k}$. This is the inverse-volume factor for two
points. With $m$ differences, each row of $AD_{\mathbf u}$ has covariance
$D_{\mathbf u}^{\mathsf T}D_{\mathbf u}$, giving the factor
$(2\pi)^{-km/2}J(\mathbf u)^{-k}$. The next lemma provides the integrability
needed to identify an exact measure on the constraint $Au_1=\cdots=Au_t$.

\begin{lemma}[Projection fibers and inverse affine volumes]
\label{lem:projection-fibers}
Let $1\le k\le d-2$ and let $P^0$ be either $\sigma_d$ or standard
Gaussian probability on $\R^d$. Every full-row-rank $k\times d$ matrix
$A$ projects $P^0$ to a Lebesgue density $p_A^0$. With
$A^\dagger=A^{\mathsf T}(AA^{\mathsf T})^{-1}$, the spherical density is
\begin{equation}\label{eq:spherical-projection-density}
 p_A^0(y)=\frac{\Gamma(d/2)}{\pi^{k/2}\Gamma((d-k)/2)}
 \frac{(1-\norm{A^\dagger y}^2)^{(d-k-2)/2}}
      {\sqrt{\det(AA^{\mathsf T})}}
 \ind_{\{\norm{A^\dagger y}<1\}}.
\end{equation}
For $\norm{A^\dagger y}<1$, its conditional law $P^0_{A,y}$ is uniform
on the sphere centered at $A^\dagger y$, of radius
$\sqrt{1-\norm{A^\dagger y}^2}$, in the affine space parallel to $\ker A$.
For the Gaussian base, $p_A^0$ is the density of $N(0,AA^{\mathsf T})$,
and the conditional point is $A^\dagger y$ plus a standard Gaussian in
$\ker A$. These conditional laws have jointly measurable versions;
unused fibers may be assigned any fixed probability law.

If $m\ge1$, $t=m+1$, and $k+m\le d-2$, then $J>0$ almost surely under
$(P^0)^{\otimes t}$ and
\begin{equation}\label{eq:inverse-affine-volume}
 \int J(\mathbf u)^{-k}\,(P^0)^{\otimes t}(d\mathbf u)<\infty.
\end{equation}
If $k+m\le d-3$, the same assertion holds with exponent $k+1$.
\end{lemma}

\begin{proof}
Split a standard Gaussian vector into the row space and kernel of $A$.
Its squared lengths in those spaces are independent chi-squares with $k$
and $d-k$ degrees of freedom. After dividing by the total length, their
ratio has the beta density. Polar integration in orthonormal row coordinates
gives \eqref{eq:spherical-projection-density}; changing those coordinates
to $y$ supplies the determinant factor. The remaining direction is
independent and uniform in the kernel. This is the spherical beta
projection law, discussed geometrically by Frankl and
Maehara~\cite{FranklMaehara1990}. Without normalization, orthogonal Gaussian
independence gives the Gaussian assertion. For measurability, generate the
kernel component by projecting an extra independent Gaussian with the Borel
matrix $I-A^\dagger A$, and normalize its length in the spherical case.
Assign fixed values on rank failures and on the null event of a zero projected
Gaussian.

For the inverse moment, first note that an affine $j$-plane $V$ satisfies
\begin{equation}\label{eq:affine-tube}
 \sigma_d\{u:\dist(u,V)\le r\}\le C^d r^{d-1-j},\qquad 0<r\le1.
\end{equation}
Indeed, cover the part of $V$ within distance one of the unit sphere by at
most $(C/r)^j$ balls of radius $r$. Each nonempty intersection with the
sphere is contained in a radius-$4r$ ball centered on the sphere. The
spherical ball bound in Lemma~\ref{lem:spherical-ball}, with the trivial
bound one when necessary, proves \eqref{eq:affine-tube}. For a standard
Gaussian point, its projection onto the orthogonal complement of the direction
space of $V$ has a translated standard Gaussian density. The density is
bounded uniformly over the translation, so the analogous estimate is
$C^d r^{d-j}$.

Gram--Schmidt gives
\[
 J(\mathbf u)=\prod_{i=2}^t
 \dist\bigl(u_i,\aff(u_1,\ldots,u_{i-1})\bigr).
\]
The predecessor affine span has dimension at most $i-2$. Thus the smallest
spherical small-distance exponent in \eqref{eq:affine-tube} is
$d-1-(t-2)=d-m>k$. Integrating its tail bounds the inverse $k$th moment
of each distance uniformly in the preceding points: for a nonnegative
distance $X$,
\[
 \E X^{-k}\le1+k\int_0^1 r^{-k-1}\Prob\{X<r\}\,dr.
\]
Integrate the last point first and repeat. This proves the finite product
moment and, since every lower-dimensional affine span has zero base mass,
full affine rank almost surely. The Gaussian exponents are larger. Replacing
$k$ by $k+1$ gives the last assertion under $k+m\le d-3$.
\end{proof}

\subsection{The exact finite measure on equal labels}

Let $f$ be a bounded nonnegative Borel function and set $\nu=fP^0$.
The measure $\nu$ is finite and need not be a probability. For a full-rank
$A$, choose the projected density and the normalized conditional probability
by
\begin{equation}\label{eq:finite-tilt}
 p_A(y)=p_A^0(y)h_A(y),\qquad
 h_A(y)=\int f\,dP^0_{A,y},\qquad
 \nu_{A,y}(du)=\frac{f(u)}{h_A(y)}P^0_{A,y}(du)
 \quad\text{when }h_A(y)>0.
\end{equation}
The conditional probability on an unused fiber may be arbitrary. In the
spherical case the displayed density is zero outside the open projection
ellipsoid. The constructions in Lemma~\ref{lem:projection-fibers} make these
versions jointly measurable.

The equality below distinguishes two natural descriptions. The left first
selects a matrix and a common label, with the label weighted by $p_A^t$,
then selects points independently on that fiber. The right first selects
independent points and then Gaussian rows annihilating their differences,
with an inverse affine-volume weight. These are finite measures, not the
two unweighted probability experiments. When $\nu$ is a probability, its
actual common label has density $p_A$, so recovering posterior replicas
will require multiplying this identity by $p_A^{1-t}$ on positive fibers.
We first identify the measures and their null sets, before that division.
The general background is coarea
\cite[Theorem 3.1]{Federer1959}; related affine simplex-volume Jacobians
occur in Drury~\cite[Lemma 1]{Drury1984}. The following proof fixes the
normalization and the versions used at the exact constrained labels.

\begin{theorem}[Finite measure for equal Gaussian labels]
\label{thm:equal-label-measure}
Let $k,m\ge1$, $t=m+1$, and $k+m\le d-2$. Let $P^0$ be either uniform
probability on $S^{d-1}$ or standard Gaussian probability on $\R^d$,
and let $\nu=fP^0$ for a bounded nonnegative Borel function $f$.
With the notation in \eqref{eq:equal-label-geometry} and
\eqref{eq:finite-tilt}, for every nonnegative Borel function
$H(A,y,\mathbf u)$,
\begin{align}
 &\int\gamma_{k,d}(dA)\int_{\R^k}p_A(y)^t
       \int H(A,y,\mathbf u)\,\nu_{A,y}^{\otimes t}(d\mathbf u)\,dy
 \notag\\
 &\quad=(2\pi)^{-km/2}\int J(\mathbf u)^{-k}\,\nu^{\otimes t}(d\mathbf u)
       \int H(A,Au_1,\mathbf u)\,\gamma_{k,d}^{F_{\mathbf u}}(dA).
 \label{eq:equal-label-measure}
\end{align}
Both sides are finite measures before the test $H$ is inserted. They are
supported on $Au_i=y$ for every $i$, full-column-rank $D_{\mathbf u}$,
and full-row-rank $A$. In the spherical case the common fiber has positive
radius. Thus \eqref{eq:equal-label-measure} is an identity of finite measures
on $(A,y,\mathbf u)$; the right side places $y$ at $Au_1$.

For a full-rank $A$, alternatively define $p_A(y)$ by the lower limit of
its radius-$1/n$ Lebesgue ball averages or by the lower limit of its Gaussian
convolutions with covariance $n^{-2}I_k$. Each choice is jointly measurable
and agrees with \eqref{eq:finite-tilt} for Lebesgue-almost every $y$.
All three versions agree at $y=Au_1$, with a finite positive value, almost
everywhere for the finite reference measure
\begin{equation}\label{eq:equal-label-reference}
 R(d\mathbf u,dA)=\nu^{\otimes t}(d\mathbf u)
                 \gamma_{k,d}^{F_{\mathbf u}}(dA).
\end{equation}
When $\nu=0$ the last assertion is vacuous. The equality is obtained using
either ball or Gaussian approximate identities in both integration orders
for the base law $P^0$, then multiplying the resulting measure identity by
the bounded density factors.
\end{theorem}

\begin{proof}
We first take $\nu=P^0$. Let $\varphi_\eta$ be either the uniform density
on the radius-$\eta$ $k$-ball or the density of $N(0,\eta^2 I_k)$. For a
continuous compactly supported test $G(\mathbf u,A)$ whose matrix support
lies in the open full-row-rank set, integrate
\begin{equation}\label{eq:equal-label-approximation}
 G(\mathbf u,A)\prod_{i=2}^t\varphi_\eta(A(u_i-u_1))
\end{equation}
against $(P^0)^{\otimes t}(d\mathbf u)\gamma_{k,d}(dA)$.

Fix $A$ first and disintegrate each point by its projected label. Given the
first label $y$, the other labels approach $y$ under either probability
approximate identity. The base density is continuous in the interior of its
support, and the conditional fiber law varies weakly continuously there.
The spherical boundary has zero projected probability. Thus the integral
converges to
\[
 \int\gamma_{k,d}(dA)\int (p_A^0(y))^t
       \int G(\mathbf u,A)(P^0_{A,y})^{\otimes t}(d\mathbf u)\,dy.
\]
To justify the passage, condition first also on $u_1$. The absolute value
of the remaining integral is at most
$\norm G_\infty(\sup p_A^0)^m$. The density suprema are uniformly bounded
on the compact matrix support: this follows from
\eqref{eq:spherical-projection-density}, whose exponent is nonnegative,
and from the Gaussian density formula. The bound is global in all the other
labels, including labels near or outside the spherical support, because
each $\varphi_\eta$ is a probability density. Dominated convergence applies
for both kernels.

In the other order, fix a tuple of full affine rank. Each row of
$AD_{\mathbf u}$ is Gaussian with covariance
$D_{\mathbf u}^{\mathsf T}D_{\mathbf u}$. Its joint density has maximum
at zero, with value
\[
 w_{k,m}(\mathbf u)=(2\pi)^{-km/2}J(\mathbf u)^{-k}.
\]
Conditional on $AD_{\mathbf u}=Z$, the matrix $A$ is a draw from
$\gamma_{k,d}^{F_{\mathbf u}}$ plus
$Z(D_{\mathbf u}^{\mathsf T}D_{\mathbf u})^{-1}D_{\mathbf u}^{\mathsf T}$.
This translation tends to zero with $Z$. The integral in
\eqref{eq:equal-label-approximation}, conditioned on the tuple, therefore
converges to
\[
 w_{k,m}(\mathbf u)\int G(\mathbf u,A)
                 \gamma_{k,d}^{F_{\mathbf u}}(dA).
\]
For either kernel its absolute value is at most
$\norm G_\infty w_{k,m}(\mathbf u)$. Indeed, the joint density of the
whole matrix $Z=AD_{\mathbf u}$ is bounded by $w_{k,m}$, while
$\prod_{i=2}^t\varphi_\eta(Z_{:,i-1})$ has Lebesgue integral one on
$\R^{k\times m}$. This argument does not require the columns of $Z$ to be
independent. The bound is integrable by Lemma~\ref{lem:projection-fibers}.
Dominated convergence gives the second integration order and equates these
two limits.

For clarity, the left measure is locally finite on the open full-row-rank
matrix set before we identify it. If a compact set in that space has matrix
projection $K_A$, its left mass is at most
\[
 \int_{K_A}\!\int (p_A^0(y))^t\,dy\,\gamma_{k,d}(dA)
 \le \sup_{A\in K_A}(\sup_y p_A^0(y))^m<\infty,
\]
because $\int p_A^0(y)\,dy=1$. Thus both sides are locally finite Borel
measures on this locally compact second-countable space; the right one is
globally finite by \eqref{eq:inverse-affine-volume}. Continuous compactly
supported tests identify them on that open set. Exhausting it shows that
the left total mass there equals the finite right total mass. The right
measure gives no mass to remaining matrices: its rows live in a space of
dimension $d-m\ge k+2$. The left measure also has no such mass because
$\gamma_{k,d}$ is full rank almost surely. Equality therefore extends to
all nonnegative Borel tests. The right measure is
supported on full affine rank by Lemma~\ref{lem:projection-fibers}, so the
same holds on the left. In the spherical case its label is in the open
ellipsoid, and its fiber radius is positive. Equivalently, a boundary fiber
could contain only one unit vector and hence could not contain a full-affine-
rank tuple with $t\ge2$.

Now multiply the base equality by $\prod_{i=1}^t f(u_i)$. The formula
\eqref{eq:finite-tilt} converts the fiber side to $p_A(y)^t\nu_{A,y}^{\otimes t}$;
the other side becomes $w_{k,m}\nu^{\otimes t}\gamma_{k,d}^{F_{\mathbf u}}$.
On a fiber with $h_A(y)=0$, the measure $fP^0_{A,y}$ is zero, so both
expressions on the fiber side are zero regardless of the unused conditional
probability. This conversion never divides on a zero fiber.
Boundedness of $f$ preserves finiteness. Inserting a test
$H(A,Au_1,\mathbf u)$ gives \eqref{eq:equal-label-measure}, since both sides
are supported on equal labels. No normalization by $\nu(\R^d)$ has been
made. In particular, the two equal-label limit calculations were made only
for the base law; no continuity of $f$ or of the reweighted fiber kernel
is used in them.

It remains to justify density evaluation on that constraint. For either
choice of $\varphi$, set
\[
 \widetilde p_A(y)=\liminf_{n\to\infty}
        \int\varphi_{1/n}(Au-y)\,\nu(du).
\]
The integral and lower limit are jointly measurable. Lebesgue
differentiation for the ball and Gaussian approximate identities makes each
version agree with \eqref{eq:finite-tilt} at almost every $y$ for each
full-rank $A$. On the left of \eqref{eq:equal-label-measure}, the common label
has a measure with Lebesgue density $p_A(y)^t$ and hence avoids the null sets
where the versions disagree or are not finite and positive. Equality of the
undivided finite measures transfers this fact to the right. There the weight
$w_{k,m}$ is finite and strictly positive at full affine rank. The reference
$R$ has finite mass $\nu(\R^d)^t$ and full affine rank almost everywhere,
because $\nu^{\otimes t}\ll(P^0)^{\otimes t}$. Restricting successively to
$\{1/j\le w_{k,m}\le j\}$ therefore transfers the same null statement to
the unweighted reference \eqref{eq:equal-label-reference}. The labels where
$p_A=0$ have zero left mass because their weight is $p_A^t$, so this
argument does not divide on them. If $\nu=0$, both finite measures and $R$
are zero. The only additional limiting assertion for a possibly
discontinuous $f$ is the ordinary almost-everywhere differentiation of its
projected $L^1$ density used in this paragraph. This proves the claimed
versions before any division by a projection density.
\end{proof}

\begin{corollary}[Posterior replica likelihood]\label{cor:replica-likelihood}
Under the hypotheses of Theorem~\ref{thm:equal-label-measure}, suppose
$p=fP^0$ is a probability. Draw $A\sim\gamma_{k,d}$ independently of
$S\sim p$, put $Y=AS$, and conditionally on $(A,Y)$ draw $t$ independent
replicas from $p_{A,Y}$. Let $Q$ be the law of their tuple together with
$A$, and let $R=p^{\otimes t}(d\mathbf u)\gamma_{k,d}^{F_{\mathbf u}}(dA)$.
Then $R$ is a probability and
\begin{equation}\label{eq:replica-likelihood}
 \frac{dQ}{dR}(\mathbf u,A)
 =(2\pi)^{-km/2}J(\mathbf u)^{-k}p_A(Au_1)^{-m}.
\end{equation}
In particular, the tuple marginal of $Q$ is absolutely continuous relative
to $p^{\otimes t}$.
\end{corollary}

\begin{proof}
Multiply the equal finite measures by $p_A(Au_1)^{-m}$ on their common
finite positive set, first with bounded truncations and then by monotone
convergence. The exponent $t=m+1$ on the left becomes one, so that side is
exactly the probability experiment defining $Q$. The right side is the
displayed density relative to $R$. Marginalization gives the last assertion.
\end{proof}

\begin{remark}[The uniform instance used with a further inverse moment]
\label{rem:extra-inverse-moment}
Suppose the replica count is $t\ge2$ and $k+t+1<d$. Then the uniform spherical
instance of Theorem~\ref{thm:equal-label-measure} applies with $m=t-1$, and
Lemma~\ref{lem:projection-fibers} also gives
$\int J^{-(k+1)}d\sigma_d^{\otimes t}<\infty$. This stronger inverse moment
controls a logarithmic volume
under a reference weighted by $J^{-k}$. Extending the likelihood argument
to an unbounded finite-entropy density requires an additional proof; it is
not included in the bounded-reweighting theorem.
\end{remark}

\subsection{Finite correlation for bounded spherical priors}

The finite-measure identity applies to both bases. The additional logarithmic
assertion needed for the first two block comparisons concerns a spherical
probability prior. It does not assert an entropy bound for every Gaussian
reweighting.

\begin{lemma}[Logarithmic integrability for spherical replicas]
\label{lem:replica-logs}
In Corollary~\ref{cor:replica-likelihood}, let $p=f\sigma_d$ with
$0\le f\le L<\infty$. Under the replica law $Q$, both $\log p_A(Y)$
and $\log J(S_1,\ldots,S_t)$ are absolutely integrable. Consequently the
tuple total correlation is finite and satisfies
\begin{equation}\label{eq:replica-correlation-upper}
 K\le -\frac{km}{2}\log(2\pi)-k\E_Q\log J-m\E_Q\log p_A(Y).
\end{equation}
\end{lemma}

\begin{proof}
Write $u_A$ for the uniform-sphere projection density. Since $p_A\le Lu_A$,
\eqref{eq:spherical-projection-density} bounds the positive logarithm of
$p_A(Y)$ by a constant depending on $d,k,L$ plus
$\tfrac12|\log\det(AA^{\mathsf T})|$. Gaussian row Gram--Schmidt expresses
that determinant as a product of independent chi-squares with degrees
$d,d-1,\ldots,d-k+1$. Their absolute logarithms are integrable, directly
from their densities at zero and infinity. For the negative part, the
support of $p_A$ lies in the ball of radius $\norm A_{\op}$ and
$-z\log z\le1/e$ for $0<z\le1$. The expected volume of that ball is
finite because a finite Gaussian matrix has all positive norm moments.

For a simplex height, condition on $(A,Y)$ and let $\rho$ be the fiber
radius. Set $q=d-k$. After translating and rescaling to the unit sphere in
$\ker A$, the posterior has density at most
\[
 L' =\frac{L u_A(Y)}{p_A(Y)}\ge1.
\]
For $a_i=\dist(S_i,\aff(S_1,\ldots,S_{i-1}))$, the tube estimate
\eqref{eq:affine-tube}, also conditional on the predecessors, gives
\[
 \Prob\{a_i\le r\rho\mid A,Y,S_1,\ldots,S_{i-1}\}
 \le\min\{1,L'C^q r^{q-1-(i-2)}\},\qquad 0<r\le1.
\]
The exponent is positive. With $\log^-x=\max\{0,-\log x\}$, tail
integration yields
\[
 \E[\log^-(a_i/\rho)\mid A,Y,S_1,\ldots,S_{i-1}]
 \le\log L'+Cq+1.
\]
Moreover the bounded density ratio makes its relative entropy finite and
\begin{equation}\label{eq:posterior-tilt-budget}
 0\le\E\log L'=\log L-\E_A\KL(p_A\Vert u_A)\le\log L.
\end{equation}
For completeness, the negative part of the density-ratio logarithm is
integrable because $-z\log z\le1/e$ against the probability density $u_A$;
the positive part is at most $\log L$.

Each pair $(A,S_i)$ has the original independent matrix--signal law.
For every fixed unit signal, rotational invariance of $A$ makes $\rho^2$
a beta variable with parameters $(d-k)/2,k/2$, whose logarithm is integrable.
Thus $\E[-\log\rho]<\infty$. Each unscaled height is at most two, so its
positive logarithm is bounded. Summing its two logarithmic parts over the
Gram--Schmidt product proves absolute integrability of $\log J$.
We may now integrate the logarithm of \eqref{eq:replica-likelihood}; this is
the finite divergence $\KL(Q\Vert R)$. Forgetting $A$ sends $R$ to
$p^{\otimes t}$, so data processing proves
\eqref{eq:replica-correlation-upper}.
\end{proof}
